\documentclass[11pt]{article}

%% ── packages ──────────────────────────────────────────────────────
\usepackage[margin=1in]{geometry}
\usepackage{amsmath,amssymb,amsthm}
\usepackage{mathtools}
\usepackage{stmaryrd}
\usepackage{mathrsfs}
\usepackage{enumitem}
\usepackage{xcolor}
\usepackage{microtype}
\usepackage[colorlinks=true,linkcolor=NavyBlue,citecolor=NavyBlue,urlcolor=SkyBlue]{hyperref}

%% ── F1R3FLY brand colours ─────────────────────────────────────────
\definecolor{NavyBlue}{HTML}{1A1A2E}
\definecolor{SkyBlue}{HTML}{3FA9F5}
\definecolor{Yellow}{HTML}{F3D630}
\definecolor{Sage}{HTML}{8BB999}

%% ── theorem environments ──────────────────────────────────────────
\theoremstyle{definition}
\newtheorem{definition}{Definition}[section]
\newtheorem{construction}[definition]{Construction}
\newtheorem{example}[definition]{Example}
\newtheorem{remark}[definition]{Remark}

\theoremstyle{plain}
\newtheorem{theorem}[definition]{Theorem}
\newtheorem{proposition}[definition]{Proposition}
\newtheorem{lemma}[definition]{Lemma}
\newtheorem{corollary}[definition]{Corollary}

%% ── notation macros ───────────────────────────────────────────────
\newcommand{\GSLT}{\mathrm{GSLT}}
\newcommand{\iGSLT}{\mathbf{iGSLT}}
\newcommand{\ciGSLT}{\mathbf{ciGSLT}}
\newcommand{\ciGSLTtc}{\mathbf{ciGSLT}_{\mathrm{TC}}}
\newcommand{\CostC}{\mathbf{Cost}\text{-}\ciGSLT}
\newcommand{\Free}{\mathsf{Free}}
\newcommand{\Forget}{\mathsf{Forget}}
\newcommand{\Imp}{\mathcal{I}}
\newcommand{\id}{\mathrm{id}}
\newcommand{\Sig}{\mathsf{Sig}}
\newcommand{\Tok}{\mathsf{Stk}}
\newcommand{\WT}{\mathbb{T}}                  % wrapped-term sort
\newcommand{\Cost}{\mathfrak{C}}              % the cost / continuation endofunctor
\newcommand{\sgn}[2]{\{#1\}_{#2}}             % signed term  {T}_s
\newcommand{\cons}{\mathbin{:}}               % temporal cons
\newcommand{\nil}{(\,)}
\newcommand{\red}{\longrightarrow}
\newcommand{\Seq}{\equiv}
\newcommand{\hashf}{\#}
\newcommand{\quot}[1]{@#1}
\newcommand{\drop}[1]{{\ast}#1}
\newcommand{\subst}[3]{#1\{#2/#3\}}
\newcommand{\compute}{\mathsf{compute}}
\newcommand{\cf}{\mathsf{cf}}
\newcommand{\Term}{\mathcal{T}}
\newcommand{\cTerm}{\widehat{\mathcal{T}}}
\newcommand{\Kp}{K_{\mathrm{p}}}
\newcommand{\Ke}{K_{\mathrm{e}}}
\newcommand{\defeq}{\;\mathrel{:=}\;}
\newcommand{\Par}{\mathbin{|}}
\newcommand{\Card}[1]{\lvert #1 \rvert}

\title{\textbf{Continued Interactive GSLTs and the\\[2pt]
Cost Endofunctor}\\[6pt]
\large A construction one level up from cost-accounted Rholang\\[2pt]
\normalsize\textit{(revised: wrapping by construction)}}

\author{L.~G.~Meredith\thanks{\texttt{lgreg.meredith@f1r3fly.io}}
\quad\textit{with}\quad Claude (Anthropic)\\[2pt]
\small F1R3FLY.io}

\date{May 2026}

\begin{document}
\maketitle

\begin{abstract}
We lift the cost-accounting construction of the rho calculus from a
single calculus to an endofunctor on a category of interactive
graph-structured lambda theories (GSLTs). The earlier work folded
phlogiston accounting into the grammar of Rholang by signing terms and
replacing the communication rule by a family of rules that rewrite
\emph{signed} terms while consuming \emph{tokens}. Here we observe that
this is a generic procedure governed by a single shape of dynamics, the
\emph{symmetric metered cut}: an interaction constructor $K$ brings two
(co-)introductions $\Kp(I,C_{\mathrm p})$ and $\Ke(J,C_{\mathrm e})$ into
contact, and a contraction $\compute(I,J,C_{\mathrm p},C_{\mathrm e})$
produces residual continuations. An interactive GSLT whose dynamics
admits this presentation, together with a computable section of its
structural-congruence quotient and a contraction compatible with
wrapping, we call a \emph{continued} interactive GSLT. The point of the
adjective is not that such a theory is already metered or wrapped --- a
continued interactive GSLT is an ordinary, un-metered calculus --- but
that it carries exactly the structure needed for the cost construction
to \emph{build} a metered version from it.

That construction is the endofunctor $\Cost$. Given a continued
interactive GSLT, $\Cost$ produces a calculus in which the continuation
slots are \emph{wrapped terms by construction}: every continuation is a
metered thunk in the grammar of the image. The decisive design move is
that this wrapping lives in the \emph{output}, so that ``every redex sits
inside a wrapper'' is a sorting invariant of $\Cost(G)$ preserved by
reduction --- not a property that must be re-established by traversing
each contractum. The source theory's contraction need not be lifted; in
the image it already maps wrapped terms to wrapped terms, and the only
run-time interaction with tokens is \emph{unwrapping}: forcing a thunk by
spending a matching token. Metering is thereby \emph{lazy} --- work is
charged when performed, not when exposed --- and the temporal token stack
is consumed exactly in step with the reductions, realising the modulus of
the underlying cut elimination as the run length itself.

We show that $\Cost$ is an endofunctor on the category $\ciGSLT$ of
continued interactive GSLTs with bisimulation-preserving, quote-faithful
morphisms. Throughout we keep $\ciGSLT$ formally distinct from the
ambient category $\iGSLT$ of interactive GSLTs: a forgetful functor
$U : \ciGSLT \to \iGSLT$ exhibits ``continued'' as a proper structural
enrichment of ``interactive,'' so that the endofunctor's domain pins down
exactly which interactive theories admit sound, lazy cost accounting.
Finally, iterated cost accounting flattens --- nested signatures multiply
and stacks-of-stacks concatenate --- so $\Cost$ is a (non-idempotent)
monad, resolving through two adjunctions with opposite embeddings: a
structural free--forgetful adjunction that detaches the apparatus, and,
over a Turing-complete base, an internalisation adjunction that dissolves
the apparatus into the base's own computation up to weak bisimulation.
\end{abstract}

\tableofcontents
\bigskip

%% ══════════════════════════════════════════════════════════════════
\section{Introduction}
\label{sec:intro}

In previous work we showed that the cost-accounting layer of Rholang ---
the phlogiston (phlo) system --- need not be a privileged runtime
extension. By treating signatures as a species of name, token stacks as
first-class terms, and the for/send operators as ranging over
\emph{signed} terms, we folded cost accounting directly into the grammar
and operational semantics of the rho calculus. The single communication
rule \textsc{Comm} became a family of rules that gate communication on
the consumption of a matching token. Because the rho calculus is the
formal core of Rholang, this revised calculus serves as the design
specification for rearchitecting phlo accounting in Rholang~1.4.

This note takes the construction \emph{one level up}, asking what the
move requires of the underlying theory and to which theories it applies.
Two observations drive the development. First, the dynamics must be
presented as a \emph{symmetric} cut: the eliminand is not a bare term
but a co-introduction carrying its own continuation, exactly as the
sender $x!(U)$ in rho carries a payload $U$ with the same status as the
receiver's continuation $T$. Second --- and this is the heart of the
revision --- a theory presented in this form carries enough structure for
the cost endofunctor to build a metered version in which the continuation
slots are \emph{wrapped terms by construction}. We are careful to keep
the two apart: the underlying theory (a \emph{continued} interactive
GSLT) is an ordinary un-metered calculus that merely \emph{admits} the
construction; the wrapping is a feature of the \emph{image} $\Cost(G)$,
not of $G$. In that image every continuation is a metered thunk in the
grammar, so the contractum is born guarded: no traversal re-wraps it, the
contraction is not lifted, and the absence of leaks is a sorting
invariant of $\Cost(G)$ rather than a dynamic obligation. The run-time
semantics of the metered calculus only ever \emph{unwraps}, forcing a
thunk by spending a token.

The $\lambda$-calculus is the foil throughout. In the rho calculus
several distinct structures are fused into $\Par$ and $\quot{\cdot}$;
pulling them apart in $\lambda$, where they are visibly separate, exposes
the construction as a genuine endofunctor and identifies precisely what
an interactive theory must supply.

Section~\ref{sec:metered-cut} isolates the symmetric metered cut.
Section~\ref{sec:wrapped} introduces wrapping by construction, derives
no-leak as subject reduction, and contrasts lazy with eager metering.
Section~\ref{sec:two-monoids} identifies the spatial monoid $K$ and the
temporal monoid $\cons$, with the stack as a reified clock and lazily
extracted modulus. Section~\ref{sec:signatures} treats signature
construction as a section of the structural-congruence quotient, with
de~Bruijn canonicalisation supplying it for $\lambda$.
Section~\ref{sec:functor} assembles the endofunctor $\Cost$ on $\ciGSLT$.
Section~\ref{sec:adequacy} records the graded adequacy statement.
Section~\ref{sec:examples} works the two anchoring instances.
Section~\ref{sec:monad} observes that iterated cost accounting flattens,
so that $\Cost$ is a monad, and exhibits its two resolutions: a
structural free--forgetful adjunction and, over Turing-complete bases, a
computational internalisation adjunction.

%% ══════════════════════════════════════════════════════════════════
\section{The symmetric metered cut}
\label{sec:metered-cut}

\subsection{Interactive GSLTs}

A \emph{graph-structured lambda theory} is a multi-sorted term theory
$G = (\Sigma, E, R)$: a signature $\Sigma$ of sorts and operators, a set
$E$ of equations imposing a structural congruence $\Seq$, and a set $R$
of rewrite rules. It is \emph{interactive} when it carries a
distinguished interacting sort $P$ with a labelled transition system,
generated by $R$ modulo $\Seq$, on which bisimulation is the intended
equivalence. The rho calculus ($P$ = processes, $\Seq$ the
$\Par$-monoid congruence with $\alpha$, $R$ = \textsc{Comm}) and the
$\lambda$-calculus ($P$ = terms, $\Seq$ = $\alpha$, $R$ = $\beta$) are
the motivating examples.

\begin{definition}[The category $\iGSLT$]
\label{def:iGSLT}
$\iGSLT$ is the category whose objects are interactive GSLTs and whose
morphisms are \emph{theory maps} --- sort- and operator-preserving
translations that respect $\Seq$ and $R$ --- that are
\textbf{bisimulation-preserving} on the interacting sort.
\end{definition}

This is the base category of the present note. The cost construction
does \emph{not} act on all of $\iGSLT$; it acts on the objects that
carry enough additional structure to be metered, which we isolate as a
subcategory in Section~\ref{sec:functor}. Keeping the two apart is
deliberate: ``interactive'' is the ambient setting, while ``continued''
(Definition~\ref{def:ciGSLT}) is the structural condition under which
sound, lazy cost accounting exists. The distinction is what the
endofunctor's domain makes precise --- an interactive GSLT is a place
where processes interact; a continued interactive GSLT is one whose
interaction is presented as a meterable cut.

\subsection{The symmetric presentation}

The cost construction acts on dynamics presented in a particular shape.
Write the generating rule of $R$ as
%
\begin{equation}
  \frac{(C_{\mathrm p}',\,C_{\mathrm e}') \;=\;
        \compute(I,\,J,\,C_{\mathrm p},\,C_{\mathrm e})}
       {\,K\bigl(\Kp(I,\,C_{\mathrm p}),\; \Ke(J,\,C_{\mathrm e})\bigr)
        \ \red\
        K'\bigl(C_{\mathrm p}',\,C_{\mathrm e}'\bigr)\,}
  \tag{$\dagger$}
  \label{eq:metered-cut}
\end{equation}
%
distinguishing:
%
\begin{description}[leftmargin=2.2em,itemsep=2pt]
  \item[the interaction constructor $K$,] which brings two operands into
    \emph{contact}. In rho, $K = {\Par}$; in $\lambda$, $K = \mathrm{App}$.
  \item[two (co-)introductions $\Kp(I,C_{\mathrm p})$ and
    $\Ke(J,C_{\mathrm e})$,] each a constructor separating a bound or
    active part ($I$, $J$) from a \emph{continuation} ($C_{\mathrm p}$,
    $C_{\mathrm e}$). In rho, $\Kp = \mathrm{for}$ with
    $I \sim (y,x)$ and $C_{\mathrm p} = T$, while
    $\Ke = {!}$ with $J \sim x$ and $C_{\mathrm e} = U$.
  \item[the contraction $\compute$,] a partial operation that, on
    contact, transforms the two continuations into residuals
    $C_{\mathrm p}', C_{\mathrm e}'$ repackaged by $K'$. In rho and
    $\lambda$ it is (semantic) substitution.
\end{description}
%
We call \eqref{eq:metered-cut} a \emph{cut}: two introductions meet and a
contraction eliminates them. The form is \emph{symmetric} in the two
operands: the eliminand is not a bare term but a co-introduction with its
own continuation. This matches rho exactly, where $x!(U)$ was never bare
--- its payload $U$ has the same status as the receiver's continuation
$T$. The earlier asymmetric form, with a bare eliminand $E$, was an
under-elaboration; the $\lambda$-calculus recovers it as the
\emph{degenerate} case in which $\Ke$ has an empty bound part, so that
the argument is its own continuation (Section~\ref{sec:examples}).

\begin{remark}[$K$ is the only constructor that needs overloading]
The cut isolates one contact site. Every cost rule interposes a token
between the operands of $K$ and drains it on contact. There is one $K$
per cut, hence one constructor to make polymorphic. ``Meter the cut'' is
the generic content of ``place valuable friction at every interaction.''
\end{remark}

%% ══════════════════════════════════════════════════════════════════
\section{Wrapping by construction}
\label{sec:wrapped}

This section describes what the cost endofunctor \emph{builds}. Given a
continued interactive GSLT $G$ --- an un-metered calculus presented as a
symmetric cut, with a section and a wrapping-compatible contraction
(Section~\ref{sec:functor}) --- the endofunctor $\Cost$ produces a
metered calculus $\Cost(G)$ whose grammar makes every continuation a
metered thunk. The wrapping discipline below is therefore a property of
the image $\Cost(G)$; the source $G$ supplies only the structure that
makes the discipline installable.

\subsection{Continuations are metered thunks in the image}

In $\Cost(G)$ the cost endofunctor adjoins a sort $\Sig$ of signatures
(Section~\ref{sec:signatures}), a \emph{wrapped-term} sort $\WT$ with the
wrapper $\sgn{\cdot}{s} : P \times \Sig \to \WT$, and a sort $\Tok$ of
token stacks
%
\begin{equation}
  S \;::=\; \nil \;\mid\; s \cons S, \qquad s \in \Sig .
  \label{eq:stack}
\end{equation}
%
The pivotal discipline is grammatical, and it is imposed by $\Cost$ on
its image: \textbf{in $\Cost(G)$ every continuation slot of $\Kp$ and
$\Ke$ has sort $\WT$.} Structural composition (parallel composition in
rho) lives at the $\WT$ level, between wrappers, so that each
redex-bearing thread is individually wrapped. Thus a continuation in the
metered calculus is not a bare process awaiting metering; it is a
\emph{metered thunk} --- $\sgn{P}{s}$ --- whose activation is gated by a
token keyed to $s$, and whose own continuations are again thunks. (In
the source $G$ these slots carried ordinary terms; re-sorting them to
$\WT$ is part of what $\Cost$ does, Construction~\ref{con:object}.)
Reading $C_{\mathrm p} = \sgn{P}{s_{\mathrm p}}$ and
$C_{\mathrm e} = \sgn{Q}{s_{\mathrm e}}$, the contraction of the image has
the type
%
\[
  \compute \;:\; \mathit{Bnd}\times\mathit{Bnd}\times\WT\times\WT
            \;\longrightarrow\; \WT\times\WT,
\]
%
mapping wrapped continuations to wrapped continuations. It need not know
that wrapping exists: substituting a wrapped term into a position yields a
wrapped term, so $\compute$ preserves the sort $\WT$ for free.

\subsection{The gated rule: force by unwrapping}

A wrapped redex is forced by consuming a matching token. The
single-signature rule is
%
\begin{equation}
  \frac{(C_{\mathrm p}',C_{\mathrm e}') = \compute(I,J,C_{\mathrm p},C_{\mathrm e})}
       {\,K\!\bigl(\sgn{K(\Kp(I,C_{\mathrm p}),\,\Ke(J,C_{\mathrm e}))}{s},
        \; s \cons S\bigr)
        \ \red\
        K\!\bigl(K'(C_{\mathrm p}',C_{\mathrm e}'),\; S\bigr)\,}
  \tag{R1}
  \label{eq:R1}
\end{equation}
%
in which the outer polymorphic $K$ combines the wrapped redex with the
token stack, the head cell $s$ is consumed (the act of \emph{unwrapping}
the redex), and the result $K'(C_{\mathrm p}',C_{\mathrm e}')$ is --- by
the typing of $\compute$ --- already built from wrapped terms. There is
no re-wrapping step and no lifted contraction. The remaining
token-distribution cases vary how the redex is split across $K$ and how
tokens are presented:
%
\begin{align}
  &K\!\bigl(K(\sgn{\Kp(I,C_{\mathrm p})}{s_1},\,\sgn{\Ke(J,C_{\mathrm e})}{s_2}),\;
        s_1 \cons s_2 \cons S\bigr)
   \ \red\
   K\!\bigl(K'(C_{\mathrm p}',C_{\mathrm e}'),\; S\bigr),
   \tag{R2}\label{eq:R2}\\[3pt]
  &K\!\bigl(K(\sgn{\Kp(I,C_{\mathrm p})}{s_1},\,\sgn{\Ke(J,C_{\mathrm e})}{s_2}),\;
        K(s_1\cons S_1,\, s_2\cons S_2)\bigr)
   \ \red\
   K\!\bigl(K'(C_{\mathrm p}',C_{\mathrm e}'),\; K(S_1,S_2)\bigr),
   \tag{R3}\label{eq:R3}
\end{align}
%
each under the same premise. Here the continuations $C_{\mathrm p},
C_{\mathrm e}$ are themselves wrapped (suppressed inner signatures), so
the residuals $C_{\mathrm p}', C_{\mathrm e}'$ remain wrapped and the next
layer of redexes is guarded without any action by the rule.

\subsection{No leak, as a sorting invariant}

\begin{definition}[Well-wrapped configuration]
A configuration of $\Cost(G)$ is \emph{well-wrapped} when it is
well-sorted in the grammar of Construction~\ref{con:object} --- in
particular, every continuation occupies a slot of sort $\WT$, so every
redex occurrence lies inside a wrapper $\sgn{\cdot}{s}$.
\end{definition}

\begin{lemma}[Subject reduction for wrapping]
\label{lem:subject}
Well-wrapping is preserved by \eqref{eq:R1}--\eqref{eq:R3}.
\end{lemma}

\begin{proof}[Proof sketch]
Each rule replaces a wrapped redex by $K'(C_{\mathrm p}',C_{\mathrm e}')$
with $(C_{\mathrm p}',C_{\mathrm e}') = \compute(\dots)$. Since
$\compute : \cdots \to \WT\times\WT$ and $K'$ packages terms of sort
$\WT$, the residual is well-sorted; the consumed cell leaves a
well-sorted stack. Redexes elsewhere are untouched. Hence the result is
well-wrapped.
\end{proof}

\begin{corollary}[No leak, by construction]
\label{cor:noleak}
In a well-wrapped configuration no redex can fire without being unwrapped,
i.e.\ without consuming a token. There is no dynamic wrapping operation;
the cost-accounting steps only ever unwrap.
\end{corollary}

The contrast with the earlier formulation is the substance of this
revision.

\begin{remark}[Eager vs.\ lazy metering]
\label{rem:lazy}
An alternative discipline re-wraps the contractum at contraction time,
traversing it to sign each exposed redex --- charging \emph{eagerly} for
every redex the step exposes. Wrapping by construction charges
\emph{lazily}: each redex is born wrapped and is charged only when
actually forced. Lazy metering pays for work performed, not work merely
exposed; redexes beneath a discarded binder are never charged. The token
stack drains exactly in step with the reductions performed, which is also
the operational reality of phlogiston consumption.
\end{remark}

\begin{remark}[Duplication needs no fresh signatures]
\label{rem:dup}
Because multiplicity is carried by the \emph{stack} (linear token
consumption) and not by key distinctness, copying a wrapped term
$\sgn{Q}{s}$ copies its key $s$ but not its tokens. Each copy must be
forced separately and so funded separately from the stack: $n$ copies of
an $s$-redex require $n$ cells keyed to $s$. Shared keys with a linear
token supply already make duplication cost, so no fresh-per-copy minting
is required. Created redexes --- where a substituted wrapped term lands in
head position --- are guarded for the same reason: the substituted
material carries its own wrapper. Thus a single mechanism (wrapped
continuations $+$ linear stack) handles duplicated and created redexes
alike, where the eager formulation needed a separate fresh-signature
discipline.
\end{remark}

%% ══════════════════════════════════════════════════════════════════
\section{Two monoids: space and time}
\label{sec:two-monoids}

The construction places two combination operators side by side, of
different character. The interaction constructor $K$ is \emph{spatial}:
it records what is in contact with what, and may carry equations ---
associative--commutative in rho (an unordered bag of processes), free in
$\lambda$ (a positional application spine). The cons operator $\cons$ of
\eqref{eq:stack} is \emph{temporal}: it records what is spent next, then
next. It is a free monoid (a list), \emph{never} commutative, because
reduction is sequential even when interaction is parallel. Each forcing
pops a cell; the order of popping is the order of steps. The stack is the
time axis of the computation reified as a term: $\Seq$ leaves it
invariant (spatial congruence consumes no time), $\red$ pops it (a step
is a tick), and its ordering is the arrow of the computation.

\begin{proposition}[Stack consumption is the modulus, lazily realised]
\label{prop:modulus}
For a terminating reduction in $\Cost(G)$, the number of cells consumed
equals the number of forced redexes, which equals the number of cuts
eliminated --- including those introduced by duplication, each charged
when forced. The consumed prefix of the temporal stack is therefore an
exact operational modulus for the cut elimination, realised by running
the reduction rather than by a separate traversal.
\end{proposition}

Where eager metering pre-computes a bound by traversing each contractum,
wrapping by construction lets the reduction \emph{be} the bound
extraction: the consumed stack is the realised cost, tight by
construction and lazy. Evaluation is cut elimination; the token stack is
the extracted modulus; metering at the granularity of forcing makes the
two coincide.

%% ══════════════════════════════════════════════════════════════════
\section{Signature construction as a section of the quotient}
\label{sec:signatures}

Tokenisation derives its expressiveness from minting signatures from
runtime data, so the construction adjoins a constructor turning terms
into signatures. We isolate what the base theory must provide.

\subsection{Quotient and section}

Let $\Term$ be the free (pre-congruence) syntax of the interacting sort
and $\cTerm = \Term/{\Seq}$ its quotient by structural congruence. Two
maps relate them, differing precisely in their relationship to $\Seq$:
%
\begin{itemize}[leftmargin=1.6em,itemsep=2pt]
  \item the \textbf{quotient} $\Term \twoheadrightarrow \cTerm$,
    \emph{respecting} $\Seq$, invertible up to congruence --- the
    structure a name constructor exposes when channels are taken up to
    $\Seq$ (in rho, $\quot{P}$ with $\quot{P}=\quot{Q}\iff P\Seq Q$). It
    is not a constructor the cost functor requires; it is the congruence
    the cut is gated modulo.
  \item a \textbf{section} $\cf : \cTerm \to \Term$, a choice of
    canonical representative per class. Composing with a
    collision-resistant digest gives the signature constructor
    $\hashf = \mathrm{digest}\circ\cf : \cTerm \to \Sig$, a one-way
    commitment to a representative. It does \emph{not} respect $\Seq$ ---
    and must not, since a commitment identifying congruent-but-distinct
    representatives would be no commitment.
\end{itemize}
%
$\Sig$ carries a commutative monoid $(\Sig,*,\nil)$ compounding
signatures, matching the compound-signature gating of \eqref{eq:R2},
\eqref{eq:R3}. Signatures never reduce: they are a rewrite-free
sub-theory meeting the behavioural theory only at the matching guard.
Hence the $\Seq$-incoherence of $\hashf$ is harmless for bisimulation,
which factors as ``$P$-bisimulation gated by signature-key matching.''

\begin{remark}[The two axes, disentangled]
In rho the section and the communication-quotient are the \emph{same}
operator $\quot{\cdot}$ viewed two ways, which is why the cost story there
appeared to need ``two axes of reflection.'' They are not two signature
schemes. One axis ($\quot{\cdot}$ up to $\Seq$) is the communication
quotient the LTS already carries; the other
($\hashf=\mathrm{digest}\circ\cf$) is the authentication section the cost
functor adds. The signature is $\hashf$ alone; the functor needs both
because metering tracks \emph{what a process does} (quotient) and
\emph{who committed to it} (section).
\end{remark}

\subsection{The \texorpdfstring{$\lambda$}{lambda}-calculus supplies a
section without a channel constructor}

The $\lambda$-calculus has no analogue of $\quot{\cdot}$. This is the
sharp test, and it passes, because the functor never required a
\emph{channel} constructor; it required a \emph{section}. For $\lambda$
the congruence is $\alpha$ and a section is a canonical representative of
each $\alpha$-class: the de~Bruijn encoding. Thus
$\hashf_\lambda = \mathrm{digest}\circ(\text{de~Bruijn})$, and a signature
commits to the de~Bruijn form. The communication-quotient role that
$\quot{\cdot}$ played in rho is, in $\lambda$, filled not by a constructor
but by $\alpha$-equality of redexes in the guard. Rho fused section and
quotient; $\lambda$ pulls them apart, and the construction is identical.

\begin{definition}[Section-equipped]
\label{def:section}
An interactive GSLT is \emph{section-equipped} when it comes with a
computable section $\cf : \cTerm \to \Term$ of its $\Seq$-quotient
(equivalently, a computable canonical-form function on the interacting
sort).
\end{definition}

%% ══════════════════════════════════════════════════════════════════
\section{The cost endofunctor}
\label{sec:functor}

\subsection{Continued interactive GSLTs}

\begin{definition}[Continued interactive GSLT]
\label{def:ciGSLT}
A \emph{continued interactive GSLT} is an interactive GSLT (an object of
$\iGSLT$) \emph{equipped with} the following additional structure:
\begin{enumerate}[label=(\roman*),itemsep=1pt]
  \item a presentation of its dynamics in symmetric metered-cut form
    \eqref{eq:metered-cut}, naming the constructors $K, \Kp, \Ke, K'$
    and the partial contraction $\compute$;
  \item a section (Definition~\ref{def:section}), i.e.\ it is
    section-equipped; and
  \item \emph{wrappability}: the contraction $\compute$ is well-sorted as
    an operation on wrapped continuations,
    $\compute : \cdots \to \WT\times\WT$ --- equivalently, $\compute$ is
    definable on metered thunks, mapping wrapped inputs to wrapped
    outputs.
\end{enumerate}
We write $\ciGSLT$ for the category whose objects are continued
interactive GSLTs and whose morphisms are theory maps that are
\textbf{bisimulation-preserving} on the interacting sort and
\textbf{quote-faithful}: injective on the reachable signature keys.
\end{definition}

The adjective \emph{continued} thus names a strict enrichment of
\emph{interactive}: clauses (i)--(iii) are data and conditions added on
top of an $\iGSLT$-object, not a restriction of it. Forgetting them
recovers the underlying interactive GSLT.

\begin{definition}[The forgetful functor]
\label{def:forget}
Let $U : \ciGSLT \to \iGSLT$ send a continued interactive GSLT to its
underlying interactive GSLT --- discarding the metered-cut presentation,
the section, and the wrappability witness --- and act as the identity on
the underlying theory map of a morphism.
\end{definition}

\begin{proposition}[The distinction is proper]
\label{prop:proper}
$U$ is faithful but neither full nor essentially surjective. Hence
$\ciGSLT$ is a genuine, proper subcategory of $\iGSLT$: strictly more
structure on objects, strictly fewer morphisms between them.
\end{proposition}

\begin{proof}[Discussion]
\emph{Faithful, not full.} $U$ forgets no information about a morphism's
action, so it is faithful. It is not full because a $\ciGSLT$-morphism
must additionally be quote-faithful: an $\iGSLT$-morphism that is
bisimulation-preserving yet collapses two distinct signature keys is a
map between the underlying objects with no lift to $\ciGSLT$. (This
refines the naive expectation that the inclusion be full: the section
structure on objects induces a real constraint on morphisms, so the
inclusion is faithful-not-full rather than full.)

\emph{Not essentially surjective.} Not every interactive GSLT admits the
added structure. Clause (i) requires the dynamics to factor as a
symmetric cut --- contact $K$ then contraction $\compute$ --- which a
theory with non-local or unfactorable rewrites need not; clause (ii)
requires a \emph{computable} section of the $\Seq$-quotient, which a
theory with undecidable structural congruence lacks; clause (iii)
requires the contraction to preserve wrapping. An interactive GSLT
failing any of these is an object of $\iGSLT$ with no preimage under
$U$.
\end{proof}

\begin{remark}[$\lambda$ exhibits the distinction]
\label{rem:lambda-distinction}
The $\lambda$-calculus is an object of $\iGSLT$ \emph{as such}: terms,
$\alpha$, $\beta$. It becomes an object of $\ciGSLT$ only once one
\emph{chooses} the added structure --- the de~Bruijn section for (ii),
the degenerate-$\Ke$ symmetric-cut presentation of $\beta$ for (i),
substitution-on-thunks for (iii) (Section~\ref{sec:examples}). The same
underlying calculus thus sits at both levels, witnessing that
``continued'' is added structure and not a property read off the
interactive GSLT alone. This is exactly the role $\lambda$ plays as a
control: each clause is a real question about it, with a real answer.
\end{remark}

The morphism conditions are dual and jointly necessary.
Bisimulation-preservation forbids distinguishing too much
(behaviourally); quote-faithfulness forbids collapsing too much
(representationally). A morphism merging two signatures would let a term
signed by one drain a token minted for the other, so the target would
gain transitions the source lacked. $\Cost$ lives exactly where both are
jointly satisfiable.

\begin{remark}[Wrappability replaces traversability]
The condition (iii) is the wrapped-by-construction successor of the
earlier ``traversable contractum'' requirement. We no longer need to walk
a contractum to re-sign its exposed redexes; we need only that the
contraction be a well-sorted operation on thunks. An opaque contraction
that destroyed wrapping --- returning bare redexes --- would violate
(iii) and is excluded, which is the same exclusion as a contraction whose
cost could not be accounted.
\end{remark}

\subsection{The functor on objects}

\begin{construction}[$\Cost$ on objects]
\label{con:object}
For $G = (\Sigma,E,R,K,\Kp,\Ke,K',\compute,\cf) \in \ciGSLT$, let
$\Cost(G)$ be the continued interactive GSLT with:
\begin{itemize}[leftmargin=1.6em,itemsep=2pt]
  \item \textbf{signature} $\Sigma$ extended by $\Sig$ (with
    $\hashf=\mathrm{digest}\circ\cf$, monoid $*,\nil$); the wrapped-term
    sort $\WT$ with $\sgn{\cdot}{s}$; the token-stack sort $\Tok$ with
    $\cons$; the polymorphic typing of $K$ over wrapped-term/stack and
    stack/stack pairs; and, crucially, the re-sorting of every
    continuation slot of $\Kp,\Ke$ to $\WT$;
  \item \textbf{equations} $E$ with the commutative-monoid equations on
    $(\Sig,*,\nil)$ and the inertness of $\Sig$ and $\Tok$ under $\red$;
  \item \textbf{rewrites} the gated family
    \eqref{eq:R1}--\eqref{eq:R3}, with the \emph{unchanged} contraction
    $\compute$ (now read at sort $\WT\times\WT$);
  \item \textbf{section} the evident lift of $\cf$ ignoring inert sorts.
\end{itemize}
\end{construction}

\begin{proposition}[Closure]
\label{prop:closure}
$\Cost(G)$ is again a continued interactive GSLT.
\end{proposition}

\begin{proof}[Proof sketch]
The gated rules are themselves in symmetric metered-cut form: the outer
polymorphic $K$ is contact, the signed redex and the token stack are the
two operands, and $\compute$ (with the stack tail) is the contraction.
$\Cost(G)$ is section-equipped by the lifted $\cf$, and wrappable because
$\compute$ preserves $\WT$ by Lemma~\ref{lem:subject}. All three clauses
of Definition~\ref{def:ciGSLT} hold.
\end{proof}

\subsection{The functor on morphisms}

\begin{construction}[$\Cost$ on morphisms]
For $f : G \to H$ in $\ciGSLT$, let $\Cost(f)$ act as $f$ on the
interacting sort and transport the adjoined structure along $f$:
signatures by $\hashf_H\circ f$ on canonical forms, wrappers and stacks
componentwise, the polymorphic $K$ structurally.
\end{construction}

\begin{theorem}[$\Cost$ is an endofunctor on $\ciGSLT$]
\label{thm:functor}
$\Cost : \ciGSLT \to \ciGSLT$ preserves identities and composition, and
$\Cost(f)$ is a $\ciGSLT$-morphism whenever $f$ is.
\end{theorem}

\begin{proof}[Proof sketch]
Object closure is Proposition~\ref{prop:closure}; identity and
composition preservation are immediate from the componentwise definition.
By the factoring of bisimulation as ``$P$-bisimulation gated by key
matching,'' and since $f$ is $P$-bisimulation-preserving and
quote-faithful, $\Cost(f)$ preserves a gated transition iff $f$ preserves
the underlying $P$-transition and respects the key match --- both hold.
Quote-faithfulness transports along the injective $\hashf_H$.
\end{proof}

%% ══════════════════════════════════════════════════════════════════
\section{Graded adequacy}
\label{sec:adequacy}

The token stack grades transitions of $\Cost(G)$ by the signature monoid
$(\Sig,*,\nil)$: a forced step is labelled by the signature(s) it
consumes. Applying OSLF --- which auto-generates a Hennessy--Milner logic
from a GSLT with an adequacy theorem --- to the graded transition system
of $\Cost(G)$ yields a \emph{graded} Hennessy--Milner logic with
modalities $\langle a\rangle_s$ indexed by the consumed signature.

\begin{theorem}[Graded adequacy, schematic]
\label{thm:adequacy}
For $\Cost(G)$ the graded Hennessy--Milner logic generated by OSLF is
adequate for quote-faithful bisimulation: two configurations satisfy the
same graded-HML formulae iff they are quote-faithfully bisimilar.
\end{theorem}

This also explains the throughput improvement mechanically. In the
translation presentation, fuel acquisition was an encoded multi-step
protocol on auxiliary channels; under $\Cost$ it is a single primitive
graded transition --- the grading absorbs the sub-protocol into the step
relation (\emph{protocol fusion}). With wrapping by construction the
fusion is total: there is no re-wrapping sub-step either, only forcing.
Together with Proposition~\ref{prop:modulus}, phlogiston is the conserved
grade, invariant under $\Seq$ (signatures inert), strictly consumed along
$\red$ (forcing drains the temporal stack), in the order recorded by
$\cons$.

%% ══════════════════════════════════════════════════════════════════
\section{Two anchoring instances}
\label{sec:examples}

\subsection{Communication (rho): the symmetric case}

Instantiate at the rho calculus: $K = {\Par}$ (associative--commutative);
$\Kp = \mathrm{for}$ with $I\sim(y,x)$ and continuation
$C_{\mathrm p} = T = \sgn{P}{s_{\mathrm p}}$;
$\Ke = {!}$ with $J\sim x$ and continuation
$C_{\mathrm e} = U = \sgn{Q}{s_{\mathrm e}}$;
$\compute$ semantic substitution; $\cf$ a normal form for the
$\Par$-congruence. Both operands are genuine introductions carrying
wrapped continuations --- this is the symmetric cut in full. Rule
\eqref{eq:R1} reads
%
\[
  \sgn{\mathrm{for}(y\!\leftarrow\! x)\,T \,\Par\, x!(U)}{s}
  \;\Par\; s\cons S
  \ \red\
  \subst{T}{\quot{U}}{y} \;\Par\; S,
\]
%
and since $T = \sgn{P}{s_{\mathrm p}}$ and $U = \sgn{Q}{s_{\mathrm e}}$,
the residual $\subst{T}{\quot{U}}{y} =
\sgn{\subst{P}{\quot{\sgn{Q}{s_{\mathrm e}}}}{y}}{s_{\mathrm p}}$ is again
a wrapped thunk, to be forced later by a token keyed to $s_{\mathrm p}$.
No re-wrapping occurs; the AC structure of $\Par$ supplies the split-redex
and split/combined-token cases \eqref{eq:R2}, \eqref{eq:R3} up to $\Seq$,
recovering the earlier five-rule lattice.

\subsection{\texorpdfstring{$\beta$}{beta}-reduction: the degenerate case}

Instantiate at the $\lambda$-calculus: $K = \mathrm{App}$ (free, no
equations); $\Kp = \lambda$ with $I = x$ and wrapped body
$C_{\mathrm p} = \sgn{M}{s_{\mathrm p}}$. The argument is not a binding
co-introduction, so $\Ke$ is \emph{degenerate}: empty bound part, the
argument being its own wrapped continuation
$C_{\mathrm e} = \sgn{N}{s_{\mathrm e}}$. The contraction is substitution,
$\compute(x,-,\sgn{M}{s_{\mathrm p}},\sgn{N}{s_{\mathrm e}}) =
(\subst{\sgn{M}{s_{\mathrm p}}}{\sgn{N}{s_{\mathrm e}}}{x},\ \nil)$. Rule
\eqref{eq:R1} reads
%
\[
  \mathrm{App}\!\bigl(
    \sgn{\mathrm{App}(\lambda x.\sgn{M}{s_{\mathrm p}},\,
                      \sgn{N}{s_{\mathrm e}})}{s},
    \; s\cons S\bigr)
  \ \red\
  \mathrm{App}\!\bigl(\subst{\sgn{M}{s_{\mathrm p}}}{\sgn{N}{s_{\mathrm e}}}{x},\; S\bigr),
\]
%
with $s = \hashf_\lambda$ of the de~Bruijn form of the redex. Every copy
of $\sgn{N}{s_{\mathrm e}}$ that substitution places into $M$ carries its
wrapper, so duplicated redexes are guarded; and where $x$ stood in head
position, the landed $\sgn{N}{s_{\mathrm e}}$ guards the newly created
application --- both by construction, with no traversal
(Remark~\ref{rem:dup}). Token supply is \emph{not} function application;
the temptation to write $(\sgn{\mathrm{App}(\lambda x.M,N)}{s}\; s)$ feeds
the key as a value and gates nothing. The token sits at the contact site
and is consumed by forcing.

%% ══════════════════════════════════════════════════════════════════
\section{The cost monad and two adjunctions}
\label{sec:monad}

Three further facts organise the construction: iterated cost accounting
flattens, so $\Cost$ carries a monad structure; and that monad relates to
its base through two different adjunctions --- one structural, one
computational --- whose embeddings have opposite characteristics.

\subsection{Flattening: \texorpdfstring{$\Cost$}{Cost} is a monad}

Applying the construction twice yields $\Cost^2(G) = \Cost(\Cost(G))$: a
calculus that meters the metering of $G$. Its signatures are signatures
of signed terms, its stacks are stacks of stacks, and its wrappers nest as
$\sgn{\sgn{P}{s_2}}{s_1}$. There is a canonical way to collapse the two
layers into one, and --- tellingly --- it is assembled entirely from the
two monoids already present in the construction.

\begin{construction}[Unit and multiplication]
\label{con:monad}
Define the \emph{unit} $\eta_G : G \to \Cost(G)$ to wrap each term with
the unit signature and present it against the freely available unit
token,
\[
  \eta_G(P) \;=\; \sgn{P}{\nil}.
\]
Since $\nil$ is the unit of $(\Sig,*,\nil)$, an $\nil$-keyed redex fires
without net resource; thus $\eta_G$ embeds $G$ as the \emph{cost-free
fragment} of $\Cost(G)$ --- the metered calculus restricted to
trivially-funded interactions. (This is consistent with the source/image
discipline of Section~\ref{sec:wrapped}: $\eta$ is precisely the wrapping
that installs apparatus at zero charge.)

Define the \emph{multiplication} $\mu_G : \Cost^2(G) \to \Cost(G)$ to
flatten the two layers by
\begin{enumerate}[label=(\roman*),itemsep=1pt]
  \item \textbf{multiplying nested signatures} in $\Sig$,
    $\sgn{\sgn{P}{s_2}}{s_1} \mapsto \sgn{P}{\,s_1 * s_2\,}$; and
  \item \textbf{concatenating the stack-of-stacks} into a single stack ---
    the free-monoid (list) multiplication on $\Tok$.
\end{enumerate}
The outer account is thereby merged into the inner: the grade of the
flattened redex is the product of the two grades, and its temporal stack
is the concatenation, in the order the two layers would have drained them.
\end{construction}

\begin{proposition}[The cost monad]
\label{prop:monad}
$(\Cost, \eta, \mu)$ is a monad on $\ciGSLT$. Its laws descend from the
laws of the two constituent monoids: the unit laws from
$s * \nil = s = \nil * s$ together with the singleton/empty laws of the
list monad, and associativity from associativity of $*$ and of list
concatenation.
\end{proposition}

\begin{proof}[Proof sketch]
$\eta$ and $\mu$ are $\ciGSLT$-morphisms. Each is quote-faithful, since
$*$ is injective in each argument up to the collision-resistance of the
digest, and bisimulation-preserving, since flattening neither creates nor
destroys gated transitions --- it relabels each by the product grade and
consolidates the (spatially commutative, temporally sequential)
consumption order consistently. Naturality is the componentwise action on
the adjoined sorts. The monad identities
$\mu \circ \Cost\eta = \id = \mu \circ \eta_{\Cost}$ reduce to the unit
laws of $*$ and of concatenation; $\mu \circ \Cost\mu = \mu \circ
\mu_{\Cost}$ reduces to their associativity.
\end{proof}

\begin{remark}[A genuine, non-idempotent resource monad]
$\Cost$ is not idempotent: $\mu_G$ is no isomorphism, because
concatenating two stacks forgets the boundary between them and
multiplying two grades forgets the factorisation. Flattening is a true
\emph{merge} of two resource accounts, not the collapse of a redundant
copy. This is the expected signature of a resource monad --- re-metering a
metered calculus yields a finer account, which $\mu$ then consolidates ---
and it distinguishes $\Cost$ from idempotent ``closure'' monads.
\end{remark}

The Eilenberg--Moore algebras of $\Cost$ are continued interactive GSLTs
equipped with a coherent \emph{discharge} map $\alpha : \Cost(A) \to A$
interpreting tokens --- a calculus that knows how to ``pay'' --- and the
Kleisli category is that of cost-accounted translations between calculi.
The structural adjunction below is the free resolution generating the
monad; the Turing-complete adjunction is a separate phenomenon.

\subsection{Adjunction I: install \texorpdfstring{$\dashv$}{-|} forget}

Let $\CostC$ be the category of \emph{cost-accounted} continued
interactive GSLTs: objects carry the installed apparatus
($\Sig$, $\Tok$, the wrapper, the gated family), morphisms preserve it.
A forgetful functor $\Forget : \CostC \to \ciGSLT$ strips the apparatus,
returning the underlying un-metered calculus; a free functor
$\Free : \ciGSLT \to \CostC$ installs it --- exactly
Construction~\ref{con:object}.

\begin{proposition}[Free--forgetful resolution]
\label{prop:adj1}
$\Free \dashv \Forget$, with unit $\eta$ of Construction~\ref{con:monad},
and the induced monad $\Forget \circ \Free$ on $\ciGSLT$ is $\Cost$.
\end{proposition}

The embedding is \emph{structural and strict}: $\Free$ adjoins sorts and
rules, $\Forget$ erases them on the nose. Crucially, forgetting is
\emph{not} behaviour-preserving in the metering sense --- erasing the
apparatus removes the gating, so $\Forget\,\Free\,G$ runs $G$'s
interactions \emph{without friction}. The structural resolution answers
``what apparatus was added?'' and lets us remove it; it does not claim the
metered and un-metered dynamics agree. Its character is
\textbf{structure-preserving, behaviour-altering}.

\subsection{Adjunction II: internalise \texorpdfstring{$\dashv$}{-|} include, over a universal base}

Restrict to the full subcategory $\ciGSLTtc \hookrightarrow \ciGSLT$ of
\emph{Turing-complete} continued interactive GSLTs. For such a $G$, the
entire metering apparatus of $\Cost(G)$ --- signature equality, stack
maintenance, gating --- is computable, hence \emph{encodable in $G$
itself}. This yields an \emph{internalisation}
\[
  \Imp_G : \Cost(G) \longrightarrow G,
\]
realising the cost-accounted semantics inside the un-metered but
universal base by simulation: tokens become encoded data, the gated rules
become an interpreter loop, and each forced step is simulated by a finite
run of $G$.

\begin{proposition}[Internalisation as an adjoint retraction]
\label{prop:adj2}
For $G \in \ciGSLTtc$ the internalisation satisfies
$\Imp_G \circ \eta_G \approx \id_G$ up to weak bisimulation, exhibiting
$\Cost(G)$ as a behavioural retract of $G$. In the bicategory of
continued interactive GSLTs, simulations, and simulation transformations,
$\eta_G \dashv \Imp_G$ is an adjoint retraction: the unit
$\eta_G$ is an (iso-up-to-bisimulation) section and the counit is the
collapsing simulation $\eta_G \circ \Imp_G \Rightarrow \id_{\Cost(G)}$.
\end{proposition}

\begin{proof}[Proof sketch]
Turing completeness of $G$ gives a computable encoding of the finite data
and decidable guards of $\Cost(G)$; standard interpreter constructions
make $\Imp_G$ a weak simulation that reflects every gated transition.
Composing with $\eta_G$ (which installs apparatus at unit cost) returns
$G$'s dynamics up to the administrative reductions of the interpreter,
i.e.\ up to weak bisimulation. The triangle identities hold up to the
2-cells witnessing these weak bisimulations, which is the content of an
adjunction internal to the simulation bicategory.
\end{proof}

The embedding is the mirror image of Adjunction~I. It is
\emph{behaviour-preserving} --- the simulation faithfully reproduces every
gated step --- but \emph{structure-encoding rather than
structure-preserving}: the clean sorts $\Sig, \Tok$ dissolve into encoded
data, and the agreement holds only up to weak bisimulation, since the
interpreter interposes administrative reductions. Its character is
\textbf{behaviour-preserving, structure-dissolving}, exactly opposite to
Adjunction~I.

\begin{remark}[Two senses of conservativity, and where they hold]
The two adjunctions express two senses in which cost accounting is
conservative over its base. Structurally (Adjunction~I, available for
every object) the apparatus is a detachable layer. Computationally
(Adjunction~II) the apparatus is already expressible, so metering is a
definitional extension up to weak bisimulation --- but this holds
\emph{only} when the base is Turing complete. The second is therefore a
genuine property of the object: universal interactive theories are
exactly those whose cost-accounted versions fold back into them, a
distinction invisible to the purely structural resolution.
\end{remark}

The cost-accounting move of the rho calculus is an instance of a generic
endofunctor on the category $\ciGSLT$ of \emph{continued} interactive
GSLTs: interactive theories presented as a symmetric metered cut
(contact $K$, two (co-)introductions $\Kp(I,C_{\mathrm p})$,
$\Ke(J,C_{\mathrm e})$, contraction $\compute$), equipped with a
computable section of their structural-congruence quotient, and with
continuation slots sorted as \emph{wrapped terms by construction}. The
functor $\Cost$ signs terms, adjoins a token stack, makes the contact
constructor $K$ polymorphic, and replaces the cut by a gated family. Its
soundness is a sorting invariant: every redex is born inside a wrapper,
reduction preserves wrapping (Lemma~\ref{lem:subject}), so no computation
proceeds for free and the contraction need not be lifted --- the
cost-accounting steps only ever unwrap. Metering is thereby lazy, paying
for work performed rather than exposed, and the temporal token stack is
consumed exactly in step with reduction, realising the modulus of cut
elimination as the run length. The cons operator is extension in time,
dual to the spatial extension of $K$; OSLF lifts $\Cost(G)$ to a graded
Hennessy--Milner logic with adequacy ``graded-HML equivalence equals
quote-faithful bisimulation.'' The $\lambda$-calculus, with $K$ and the
section visibly distinct and $\Ke$ degenerate, shows the construction
requires a section --- de~Bruijn canonicalisation --- not a channel
constructor; rho's apparent need for two axes of reflection was the
self-overlap of one operator serving as both quotient and section.

Finally, the construction is deliberately sited one rung below the
ambient category: $\Cost$ is an endofunctor on $\ciGSLT$, not on
$\iGSLT$, and the forgetful functor $U : \ciGSLT \to \iGSLT$
(Definition~\ref{def:forget}, Proposition~\ref{prop:proper}) records that
``continued'' is a proper structural enrichment of ``interactive.'' The
content of the note is therefore as much about \emph{where} cost
accounting lives as about what it does: an interactive GSLT is a setting
for interaction; a continued interactive GSLT is one whose interaction is
presented as a meterable cut, and it is precisely on these that the cost
endofunctor exists --- indeed, as a monad
(Section~\ref{sec:monad}), since iterated metering flattens. Its two
resolutions sharpen the picture: the structural free--forgetful
adjunction shows the apparatus is detachable, while over a Turing-complete
base the internalisation adjunction shows it is already expressible, so
that cost accounting is a definitional extension up to weak bisimulation.
That second adjunction holds only for universal bases, and so quietly
begins to sort interactive theories by a structural property --- a thread
we leave for the sequel.

\end{document}
